\scanpage{302}
\mspage{43}
\deleted{Introduisant les intervalles critiques $J,J',J''$ qui
contiennent \deleted{\unsure{J}} $u,u',u''$ respectivement, on est ramené}
\deleted{\unsure{à}} Si les flèches en question ont comme portée
intrinsèque $[i-1,i]$, $[i'-1,i']$, $[i''-1,i'']$ respectivement,
l'hypothèse est
\[
\begin{aligned}
 g(i'-1)&\leq i''-1,\quad g(i')\geq i''\\
 f(i-1)&\leq i'-1,\quad f(i)\geq i'
\end{aligned}
\]
\deleted{\unsure{et}} comme $g$ est croissant, on a d'ailleurs
\[
\begin{aligned}
 gf(i-1)&\leq g(i'-1)\ (\leq i''-1),\\
 gf(\unsure{u|i})&\geq g(i')\ (\geq i''),
\end{aligned}
\qquad\text{cqfd.}
\]

Ainsi, le cor. 1 plus haut obtient un sens précis,
et il est démontré du $=$ coup, vu que les
correspondances du cor. 1 sont données par
\[
\underset{\substack{\rotatebox{-90}{$\in$}\\
\operatorname{Hom}_{\underline{\mathrm{Ty}}}(\unsure{\tau|t},\unsure{\mathbb{\Delta}|\Delta}^{\unsure{\varepsilon|+}})}}{f}
\longmapsto\mathrm{Tr}(f)^*(\underline{u})
\]
où $\underline{u}$ est l'unique flèche de transition de $\mathbb{\Delta}^\varepsilon$.

On va utiliser ces sorites pour donner,
pour toute $\tau\in\mathrm{Ty}$, une description combinatoire
simple des catégories $\underline{\mathrm{Ty}}_{/\tau}$ et $ {}_{\tau\backslash}\underline{\mathrm{Ty}}$, ou du
moins de leurs ensembles d'objets.

\uline{1\textsuperscript{o}} \uline{Description de} $\operatorname{Ob}\,{}_{\tau\backslash}\underline{\mathrm{Ty}}\deleted{\unsure{/\tau}}$.
À $f\in\operatorname{Ob}\,{}_{\tau\backslash}\underline{\mathrm{Ty}}\deleted{\unsure{\tau}}$ i.e.
\[
 f\colon\tau\longrightarrow\tau'\qquad\text{i.e. } f\colon I_\tau\unsure{\mathbin{-}|\longrightarrow}I_{\tau'}
\]
dans $\underline{\mathrm{Ty}}$, est associé $E\colon E_f=\deleted{\unsure{\mathrm{Cr}}}\mathrm{Tr}(\tau,f)$\inserted{$\subset\mathrm{Tr}(\tau)$}, l'ens des flèches
\inserted{de transition} critiques (i.e. non dégénérées) par $f$. À chaque
$u\in\deleted{\unsure{\mathrm{Tr}}}E$, associons l'entier $\operatorname{long}(\pi(f,u))$,
% archive sheet number: 301
